\chapter{Conception de l'outil HOMO-CI}
\label{ch:conception}

Ce chapitre explique ce que l'outil doit faire et pourquoi il est
construit comme il est construit. Le code lui-même est détaillé
au chapitre suivant.

\section{À quoi sert l'outil}

HOMO-CI répond à une question simple~: \emph{comment garder le
dossier d'homologation aligné sur l'état réel du SI, sans passer
des semaines à le réécrire à la main~?}

Concrètement, il fait quatre choses à chaque exécution~:

\begin{enumerate}
  \item Il \textbf{collecte} les preuves techniques en interrogeant
    directement les sources~: API Azure pour les NSG, Trivy pour
    les CVE des images Docker, LDAP pour l'inventaire AD, API
    Wazuh pour les règles et alertes.
  \item Il \textbf{stocke} ces preuves dans une base locale
    (SQLite + journal JSONL), avec un horodatage et un hash pour
    chaque preuve.
  \item Il \textbf{compare} avec l'exécution précédente et
    identifie ce qui a changé (création, modification,
    suppression).
  \item Il \textbf{régénère} le dossier d'homologation au format
    Markdown puis PDF, en suivant la structure ANSSI, et le signe
    avec un hash SHA-256 chaîné au dossier précédent.
\end{enumerate}

\section{Cinq principes directeurs}

Les choix techniques de l'outil suivent cinq principes~:

\begin{description}
  \item[Source de vérité unique.] L'état attendu du SI est dans
    le code (scripts Terraform, manifestes Docker, configuration
    Wazuh versionnés en git). L'écart avec l'état réel collecté
    par l'outil est une \emph{dérive} qu'il faut documenter.
  \item[Idempotence.] Exécuter l'outil deux fois de suite sur le
    même état produit la même sortie. Pas d'effets de bord, pas
    de compteurs qui s'incrémentent au hasard.
  \item[Traçabilité.] Chaque preuve a un UUID, un timestamp, un
    hash de payload, une signature de source. Chaque dossier
    généré référence le hash du dossier précédent
    (\emph{hash chain}). On peut reconstruire l'historique
    complet.
  \item[Découplage par messages.] Les watchers ne se parlent
    pas entre eux. Ils écrivent dans un \emph{evidence store}
    commun (SQLite + JSONL). Si l'API Azure tombe, le watcher AD
    continue à tourner.
  \item[Le pipeline est lui-même un SI à sécuriser.] Moindre
    privilège pour les credentials Azure, secrets en variable
    d'environnement (idéalement en Key Vault), signature
    cryptographique du dossier produit.
\end{description}

\section{Architecture en six couches}

L'outil est organisé en six couches empilées~:

\begin{figure}[h]
\centering
\begin{small}
\begin{verbatim}
+----------------------------------------------+
| Couche 6 — Gouvernance (KRI, alertes)        |
+----------------------------------------------+
| Couche 5 — Publication (PDF signé + Blob)    |
+----------------------------------------------+
| Couche 4 — Génération (Jinja2 + pandoc)      |
+----------------------------------------------+
| Couche 3 — Détection de changement (Diff)    |
+----------------------------------------------+
| Couche 2 — Stockage (SQLite + JSONL)         |
+----------------------------------------------+
| Couche 1 — Collecte (watchers parallèles)    |
+----------------------------------------------+
\end{verbatim}
\end{small}
\caption{Architecture en six couches de HOMO-CI.}
\label{fig:archi-pipeline}
\end{figure}

\paragraph{Couche 1 --- Collecte.}
Un \emph{watcher} par source de données~: \texttt{wc\_nsg} pour
les NSG Azure, \texttt{wc\_cve} pour Trivy, \texttt{wc\_ad} pour
Active Directory, \texttt{wc\_siem} pour Wazuh. Chaque watcher
implémente la même interface \texttt{BaseWatcher} et produit une
liste d'\texttt{Evidence}.

\paragraph{Couche 2 --- Stockage.}
SQLite pour le requêtage rapide, JSON Lines pour les snapshots
historiques. Les deux formats sont synchronisés à chaque
collecte.

\paragraph{Couche 3 --- Diff.}
Compare deux snapshots et produit une liste de \texttt{Change}
typés (\texttt{CREATED}, \texttt{UPDATED}, \texttt{DELETED}).
Complexité $O(N+M)$ par indexation sur la clé
\texttt{(source, ref)}.

\paragraph{Couche 4 --- Génération.}
Un template Jinja2 (\texttt{dossier.md.j2}) prend les preuves
collectées et les rend en Markdown structuré selon les sections
ANSSI. Pandoc convertit ensuite en PDF, WeasyPrint applique le
style CSS.

\paragraph{Couche 5 --- Publication.}
Le PDF est signé (hash SHA-256 chaîné au précédent), poussé sur
un Blob Storage Azure, et son hash est journalisé.

\paragraph{Couche 6 --- Gouvernance.}
Indicateurs simples~: nombre de changements par cycle, fraîcheur
du dernier dossier, statut des watchers. Alerte si pas de run
depuis plus de 26~heures.

\section{Le modèle de données}

Trois entités centrales, toutes immutables et validées par
Pydantic~:

\begin{description}
  \item[\texttt{Evidence}.] Une preuve d'audit collectée. Porte
    un UUID, un timestamp UTC, une source, une catégorie, une
    référence métier (\emph{ex.} \texttt{nsg-dmz/Allow-HTTP}),
    un payload structuré et son hash SHA-256.
  \item[\texttt{Change}.] Un delta entre deux snapshots. Type
    (créé, modifié, supprimé), référence à l'évidence avant et
    après.
  \item[\texttt{DossierSnapshot}.] L'instantané du dossier
    produit. Version, date, nombre de preuves, hash du PDF,
    hash du parent.
\end{description}

\section{Cartographie des preuves vers le cadre ANSSI}

Chaque watcher déclare statiquement à quelles étapes du guide
ANSSI ses preuves contribuent. Par exemple, le watcher NSG
déclare~:

\begin{lstlisting}[language=Python,caption={Mapping NSG vers le
référentiel ANSSI / ISO.}]
class NSGWatcher(BaseWatcher):
    SCOPE = Scope(
        anssi_steps=(4, 6),       # cartographie SI, mesures
        iso_controls=("8.20", "8.21"),  # Networks security
        asvs_controls=(),
    )
\end{lstlisting}

C'est ce mapping qui permet au template Jinja2 de placer chaque
preuve dans la bonne section du dossier final.

\section{Ce que l'outil ne fait pas (encore)}

Pour cadrer~:

\begin{itemize}
  \item Il ne fait pas d'analyse de risque automatique. EBIOS RM
    reste manuel. L'outil alimente seulement les sections
    techniques du dossier.
  \item Il ne corrige rien. S'il détecte une dérive (NSG ouvert,
    CVE critique), il la documente mais ne ferme pas la règle.
    Une boucle de remédiation est une perspective future.
  \item Il n'a pas d'interface web. Tout passe par CLI
    (\texttt{homo-ci collect}, \texttt{homo-ci diff},
    \texttt{homo-ci render}, \texttt{homo-ci publish}). Pour un
    PFE, c'est suffisant~; pour un produit, il en faudrait une.
\end{itemize}
